\section{Visual and Catalog Review}

\subsection{Purpose}

The neighbor-profile analysis nominates signs that behave similarly. It does not show that the signs look alike. This milestone therefore creates a visual and catalog review pack: a prioritized list of candidate pairs, same-text examples, source-access status, and a protocol for deciding whether a pair should be merged, linked as variants, kept distinct, or held pending.

The method follows two constraints from prior scholarship. First, ICIT-style grouping uses similar sign-neighbor frequencies to identify paradigmatic behavior \citep{fulsICIT}. Second, allograph claims require image and context review: mirroring, writing direction, space-saving, regional anomalies, carving errors, and catalog errors must be checked before reducing the sign list \citep{daggumati2021allographs}. The Harappa review of Fuls' catalog emphasizes the catalog's visual documentation and sign-reference function \citep{harappaCatalog2026}, but the present workspace does not yet contain the full sign-image coverage needed for a final verdict.

\subsection{Source Status}

\begin{table}[htbp]
\centering
\begin{tabular}{p{0.24\textwidth}p{0.24\textwidth}p{0.40\textwidth}}
\toprule
\textbf{Source} & \textbf{Status} & \textbf{Use} \\
\midrule
Committed workspace image corpus & Unavailable & No committed sign or artifact image assets are present; scratch extracts remain local. \\
Fuls 2023 SearchInside local PDF & Partial local sample & Searchable 35-page sample with embedded sign-list images; includes signs \texttt{820} and \texttt{920}. \\
ICIT/Epigraphica & Catalog described, images not harvested & Primary catalog target. \\
Deutsche Digitale Bibliothek & Bibliographic record only & Confirms Fuls 2023 catalog; contents only openly accessible. \\
Harappa catalog review & Public descriptive source & Confirms catalog scope and importance. \\
Daggumati and Revesz 2021 & Open-access method source & Supplies allograph-review checks. \\
CISI volumes & Not locally available & Artifact-image authority. \\
\bottomrule
\end{tabular}
\caption{Current source status for visual and catalog review.}
\label{tab:visual-source-status}
\end{table}

\subsection{Partial Visual Evidence Found}

The new local SearchInside sample is useful but partial. It gives direct visual evidence for two signs that touch active review pairs: ICIT \texttt{820}, relevant to \texttt{817}/\texttt{820} and \texttt{820}/\texttt{861}; and ICIT \texttt{920}, relevant to \texttt{692}/\texttt{920}. This upgrades those pairs from purely distributional candidates to one-sided visual-review candidates. It does not yet license a merge or variant claim, because the counterpart signs \texttt{817}, \texttt{861}, and \texttt{692} still need matching visual evidence.

\begin{table}[htbp]
\centering
\begin{tabular}{llll}
\toprule
\textbf{Sign} & \textbf{PDF page} & \textbf{Related pair(s)} & \textbf{Status} \\
\midrule
\texttt{820} & 21 & \texttt{817}/\texttt{820}; \texttt{820}/\texttt{861} & One-sided visual evidence \\
\texttt{920} & 35 & \texttt{692}/\texttt{920} & One-sided visual evidence \\
\bottomrule
\end{tabular}
\caption{Partial visual evidence located in the local Fuls SearchInside PDF sample.}
\label{tab:partial-visual-evidence}
\end{table}

\subsection{Review Tiers}

The script \texttt{scripts/visual-catalog-review.ps1} ranks the top 80 distributional pairs into review tiers. The ranking favors high cosine similarity but penalizes sparse evidence. All rows remain marked as \emph{Pending} for visual/catalog status.

\begin{table}[htbp]
\centering
\begin{tabular}{lr}
\toprule
\textbf{Tier} & \textbf{Pairs} \\
\midrule
Tier 1 & 4 \\
Tier 2 & 3 \\
Tier 3 & 26 \\
Sparse-data hold & 47 \\
\bottomrule
\end{tabular}
\caption{Review tiers in the first visual/catalog candidate pack.}
\label{tab:visual-review-tiers}
\end{table}

\subsection{Priority Pairs}

\begin{table}[htbp]
\centering
\begin{tabular}{llrrrr}
\toprule
\textbf{Sign A} & \textbf{Sign B} & \textbf{Cosine} & \textbf{Tokens A} & \textbf{Tokens B} & \textbf{Tier} \\
\midrule
817 & 861 & 0.965020 & 221 & 272 & 1 \\
817 & 820 & 0.956926 & 221 & 244 & 1 \\
705 & 706 & 0.953396 & 219 & 102 & 1 \\
692 & 920 & 0.953170 & 77 & 153 & 1 \\
526 & 527 & 0.982841 & 13 & 63 & 2 \\
090 & 565 & 0.953096 & 195 & 15 & 2 \\
400 & 565 & 0.950984 & 502 & 15 & 2 \\
820 & 861 & 0.924203 & 244 & 272 & 3 \\
033 & 034 & 0.908290 & 523 & 184 & 3 \\
\bottomrule
\end{tabular}
\caption{Selected visual/catalog review candidates.}
\label{tab:visual-priority-pairs}
\end{table}

\subsection{Review Protocol}

Each pair must pass through six questions:

\begin{enumerate}
  \item \textbf{Shape}: Are the signs graphically similar enough to be variants?
  \item \textbf{Direction}: Could mirroring or reversal explain the difference?
  \item \textbf{Space}: Could crowding or line-end compression explain the difference?
  \item \textbf{Context}: Do positions, neighbors, sites, and artifact types support functional equivalence?
  \item \textbf{Independence}: Do the occurrences reflect independent texts rather than a repeated-object effect?
  \item \textbf{Verdict}: Merge, variant-link, keep distinct, or hold pending.
\end{enumerate}

\subsection{Outputs}

\begin{itemize}
  \item \texttt{outputs/visual\_catalog\_review\_candidates.csv}
  \item \texttt{outputs/visual\_catalog\_review\_examples.csv}
  \item \texttt{outputs/visual\_catalog\_review\_protocol.csv}
  \item \texttt{outputs/visual\_catalog\_source\_status.csv}
  \item \texttt{outputs/visual\_catalog\_found\_signs.csv}
  \item \texttt{outputs/visual\_catalog\_review.tex}
\end{itemize}

\subsection{Interpretation}

The first review pack identifies \texttt{817}/\texttt{861}, \texttt{817}/\texttt{820}, \texttt{705}/\texttt{706}, and \texttt{692}/\texttt{920} as the strongest immediate candidates for visual/catalog inspection. No sign merge is proposed at this stage.
